% Owner: ekagrahvinay-netizen
% Section body only; packages and colours are loaded by ../main.tex.
% Based on the five categories in the team's Lecture 2 environment slide.
\section{Environment Analysis}\label{sec:environment}

The Air Traffic Control Multi-Agent System (ATC-MAS) is designed for a
\emph{simulated, batch-processing} environment, not live air traffic control.
It reads a GraphML airspace network and JSON aircraft and flight-plan data,
validates them, and predicts aircraft trajectories to identify possible
separation conflicts. The model uses horizontal and vertical separation
thresholds of 5 nautical miles and 1000 feet, respectively \cite{icao4444}. We classify
its environment using the five properties discussed in Lecture 2
\cite{urvlect2,wooldridge2009}; an analogous, broader classification appears in~\cite{russell2020}.

\subsection{Environment Characteristics}

\paragraph{Accessible versus inaccessible}
Real ATC is \emph{inaccessible}: surveillance can be noisy or incomplete,
and pilot intentions cannot always be observed. Our \emph{modelled scenario}
is classified as \emph{accessible}, because the simulation receives the
scenario information required by its calculations. This does not imply
that every agent has access to all data; individual crews may receive only
the information needed to perform their roles.

\paragraph{Deterministic versus non-deterministic}
Real ATC is \emph{non-deterministic}: weather, aircraft behaviour and pilot
responses introduce uncertainty. The simulated trajectory and conflict
calculations are \emph{deterministic} for the same valid inputs and fixed
rules. Outputs proposed by language-model-based agents may still vary;
such proposals should not be confused with the deterministic safety checks.

\paragraph{Episodic versus non-episodic}
Real ATC is \emph{non-episodic}, since a clearance can affect later traffic
situations and decisions. Our simulation is \emph{episodic across runs}:
each submitted scenario is processed independently. It is
\emph{non-episodic within a run}, because conflict findings and subsequent
resolution decisions depend on earlier processing steps.

\paragraph{Static versus dynamic}
Real ATC is \emph{dynamic}, since aircraft keep moving and conditions
change while decisions are made. The simulated scenario is \emph{static}
with respect to external inputs during each batch run: no new live radar
or weather observations arrive. This does not mean the model ignores
motion; positions change along the trajectories being calculated.

\paragraph{Discrete versus continuous}
Real aircraft movement is predominantly \emph{continuous}: position,
time, altitude and speed vary continuously. Our simulation is
\emph{hybrid and discretised by design}: graph routes and a finite set of manoeuvres coexist with numerical trajectory
quantities. Candidate resolutions must still be assessed against the
model's numerical separation constraints.

\subsection{Implications for System Design}

These characteristics motivate a separation of deterministic verification
from agent coordination. Data Management validates the supplied files,
and Surveillance \& Conflict calculates potential separation violations.
Conflict-free scenarios proceed to Reporting; scenarios with conflicts
are routed to Tactical Resolution. The Safety Validation \& QA crew
independently checks suggested manoeuvres, rejecting unsafe alternatives
and allowing bounded retries before reporting the final status.

Because the environment is accessible, deterministic and static during
a run, the system can afford deliberative planning: as Lecture~2 notes,
planning with an explicit model of the world can reach optimal solutions
when the environment does not change while the plan is being
made~\cite{urvlect2}. At the same time, the environment is non-episodic
within a run: a manoeuvre that resolves one conflict can create another.
This is why every proposed manoeuvre is re-checked by the independent
Safety Validation \& QA crew before it is reported.

The resulting design supports deliberative coordination without treating
language-model output as a safety guarantee. This is an academic
simulation, not a real-world clearance system.

